\documentclass[11pt]{article}

\usepackage[T1]{fontenc}
\usepackage{amsmath,amssymb,amsthm}
\usepackage{booktabs}
\usepackage{enumitem}
\usepackage{geometry}
\usepackage[hidelinks]{hyperref}
\usepackage{microtype}

\geometry{margin=1in}
\setlist{nosep}

\newtheorem{definition}{Definition}[section]
\newtheorem{theorem}[definition]{Theorem}
\newtheorem{proposition}[definition]{Proposition}
\newtheorem{example}[definition]{Example}
\newtheorem{remark}[definition]{Remark}

\newcommand{\Build}{\mathcal{B}}
\newcommand{\Obs}{\mathcal{O}}
\newcommand{\View}{\mathcal{V}}
\newcommand{\Rel}{\mathcal{R}}
\newcommand{\Work}{\mathsf{Work}}
\newcommand{\Span}{\mathsf{Span}}

\title{AGI-Oriented Reproducible Builds as\\
Relational, Observation-Indexed GSLT Integrity}
\author{MeTTapedia Formalization Project}
\date{August 2026}

\begin{document}
\maketitle

\begin{abstract}
Masayuki Hatta proposes seven requirements for AGI-oriented reproducible
builds, extending ordinary build reproducibility to evolving and
self-modifying systems.  We give a machine-checked account in which a build is
a proof-relevant relation rather than necessarily a function, and
reproducibility is indexed by both a declared source view and an artifact
observation.  The formulation separates rebuildability, per-source agreement,
declared-input sufficiency, replayability, result-level hosting, proof-fibre
fidelity, provenance, receipt currentness, and verification cost.

The theory proves identity and composition laws for GSLT-to-GSLT
transformations, exposes hidden intermediate dependencies, and recovers
functional transport only from an explicit representability certificate.  It
instantiates Hatta's R1--R5 as a typed profile, proves R6 for a characterized
class of current protected-family-preserving self-modifications, formalizes
Wheeler's diverse double-compiling assumptions and Thompson's trusting-trust
attack, and treats R7 as coverage plus a Work--Span--storage feasibility
discipline.  Observer-indexed legibility is admitted only after exhibiting two
extensionally equal realizations with different attained verification ranks.
Dynamic-link provenance composes along relational routes but does not, without
an additional rule, determine a legal adjudication.  All positive claims are
paired with negative controls, including sampled-but-not-universal
verification and semantically equal routes with different provenance.
\end{abstract}

\section{Attribution and scope}

Hatta's seven requirements and the governance motivation come from
\cite{hatta2026,hatta2025}.  Complete input closure and purely functional
deployment are informed by Dolstra and colleagues
\cite{dolstra2006,dolstra2004}.  The distinction between task semantics,
dependency discovery, rebuilding, and scheduling follows the build-systems
analysis of Mokhov, Mitchell, and Peyton Jones \cite{mokhov2020}.  Signed
layouts and link evidence use the conceptual separation of \emph{in-toto}
\cite{intoto2019}.  Trusting trust and diverse double-compiling are due to
Thompson and Wheeler \cite{thompson1984,wheeler2009}.  The empirical and
engineering reproducible-build background is provided by
\cite{lamb2021,malka2025}.

The following are MeTTapedia extensions rather than claims attributed to those
sources:

\begin{itemize}
  \item the proof-relevant relational and observation-indexed foundation;
  \item the GSLT/GSLT-IL/OSLF composition and hosting theorems;
  \item the exact replay equivalence under stated adequacy, completeness,
        ordering, initial-state, and currentness premises;
  \item the protected-family R6 theorem for the characterized modification
        class;
  \item observer-indexed weakest-audit-view legibility and its strict
        separation witness;
  \item the R7 Work--Span--storage algebra and the dynamic-link relational
        provenance construction.
\end{itemize}

No theorem here converts Hatta's governance analogy into a legal conclusion.
No claim about arbitrary self-modifying AGI follows from the restricted R6
result.

\section{Relational build semantics}

Let $S$ be a type of complete source and environment states, $A$ a type of
artifacts, $D$ a type of declared input views, and $O$ a type of observable
artifact outcomes.  None of these types is required to be finite.

\begin{definition}[Proof-relevant build]
A build is a type-valued relation
\[
  \Build : S \to A \to \mathsf{Type}.
\]
An occurrence is a triple $(s,a,e)$ with $e : \Build(s,a)$.  Distinct evidence
terms are retained even when their source and artifact endpoints coincide.
\end{definition}

A declared input view is a map $d:S\to D$.  An artifact observation is a map
$o:A\to O$.  Identity observation gives bit- or value-exact reproducibility;
a coarser observation may express an explicitly justified tolerance or
acceptance discipline.

\begin{definition}[Separated obligations]
For a build $\Build$ and observation $o$:
\begin{align*}
\mathsf{Rebuildable}(\Build)
  &:\!\Longleftrightarrow \forall s,\; \exists a,\; \Build(s,a),\\
\mathsf{Consistent}(\Build,o)
  &:\!\Longleftrightarrow
    \forall s,a,b,\;\Build(s,a)\to\Build(s,b)\to o(a)=o(b),\\
\mathsf{Reproducible}(\Build,o)
  &:\!\Longleftrightarrow
    \mathsf{Rebuildable}(\Build)\wedge\mathsf{Consistent}(\Build,o).
\end{align*}
Declaration sufficiency is the stronger cross-source property
\[
  d(s)=d(t) \wedge \Build(s,a) \wedge \Build(t,b)
  \;\Longrightarrow\; o(a)=o(b).
\]
Declared-view reproducibility combines rebuildability with declaration
sufficiency.
\end{definition}

This separation closes several common loopholes.

\begin{proposition}[Vacuity and hidden input controls]
The formal development proves all of the following.
\begin{enumerate}
  \item Reproducibility implies rebuildability.
  \item Rebuildability alone does not imply reproducibility: one source may
        admit two observably different artifacts.
  \item Consistency alone does not imply rebuildability: an empty execution
        fibre is vacuously consistent.
  \item Declaration sufficiency implies per-source consistency.
  \item Per-source reproducibility does not imply declaration sufficiency: a
        hidden hardware bit can change the exact artifact while every complete
        source remains deterministic.
\end{enumerate}
\end{proposition}

An executable reproducer through $d$ is stronger than fibrewise sufficiency:
it computes the observed artifact from every declared value.  The converse
requires the exact reachability or inhabited-codomain premise needed to define
the reproducer away from realized declarations.  This distinction prevents a
choice function outside the image from being smuggled into the definition.

\subsection{Observation refinement}

If an exact observation factors through a coarser observation, exact
reproducibility transports downward.  Transport in the opposite direction is
not automatic.  A numerical tolerance is accepted only through a supplied
soundness law; it is not presumed transitive and therefore is not silently
treated as a quotient.

The identity observation is thus the finest ordinary case, not a synonym for
the general theory.  This is important for AGI systems whose accepted outcome
may be behavioral, probabilistic, proof-relevant, or resource-sensitive.

\section{GSLT-to-GSLT composition}

The core theory composes builds relationally.  For
$\Build_1:S\to M\to\mathsf{Type}$ and
$\Build_2:M\to A\to\mathsf{Type}$, the composite evidence retains an
intermediate $m:M$ and witnesses of both stages:
\[
  (\Build_2\circ\Build_1)(s,a)
  := \sum_{m:M}\Build_1(s,m)\times\Build_2(m,a).
\]
Consequently, repeated or alternative derivations are not collapsed merely
because they share endpoints.

\begin{theorem}[Identity and composition]
Identity builds are reproducible.  Reproducible stages compose when their
declared source views include the dependency closures required by both stages
and their intermediate observations are aligned.  Composition is associative
up to proof-fibre equivalence, rather than by erasing the intermediate
witness.
\end{theorem}

\begin{example}[Hidden intermediate input]
Two stages may each be reproducible relative to their local declarations while
the composite is not reproducible relative to a declaration that hides a
property read by the downstream stage.  The formal counterexample keeps the
middle artifact observable long enough to show exactly where sufficiency
fails.
\end{example}

Certified plans retain stage order, declared source demand, and observation
alignment.  A plan that omits a dependency is rejected even when a particular
sample execution happens to return the expected endpoint.

\section{Hosting, proof fibres, and open GSLT-IL routes}

Result-level hosting and proof-relevant hosting are different interfaces.  A
behavioral hosting certificate transports statements about produced results at
the declared observation.  It does not imply that source and target execution
fibres are equivalent.  Proof-relevant hosting supplies that stronger
equivalence explicitly.

\begin{theorem}[Observation-exact transport]
If a behavioral host preserves and reflects the selected artifact observation,
then source reproducibility transports to the hosted build at that observation.
No transport to a strictly finer observation follows without another
certificate.
\end{theorem}

The negative target called \texttt{CBad} is deterministic but invents behavior.
It demonstrates why determinism and forward preservation do not establish
no-invention or reflection.

GSLT-IL transport is relational by default.  A relation becomes a direct
function only through a proof-relevantly total and deterministic
representation certificate.  Open choice and partial routes remain valid
relations but cannot be functionalized.  Exact returned-image theorems remain
scoped to the returned image; commands outside it are retained as explicit
counterexamples.

The OSLF bridge follows the same discipline.  Syntactic translation alone does
not establish operational reproducibility.  The bridge uses a named
operational correspondence law, with both a satisfying instance and a failing
translation witness.

\section{Hatta R1--R5 as a typed profile}

The formal profile stores semantic evidence, not seven Boolean flags.  Its
fields are summarized in Table~\ref{tab:r1r5}.

\begin{table}[ht]
\centering
\small
\begin{tabular}{@{}p{0.08\linewidth}p{0.36\linewidth}p{0.45\linewidth}@{}}
\toprule
Req. & Formal object & Required evidence \\
\midrule
R1 & declared input view & declaration sufficiency at the selected observation \\
R2 & reproducibility budget & rebuildability and observational consistency \\
R3 & toolchain and hardware axes & each is declared or proved observationally abstracted \\
R4 & attestation discipline and layout & authorized, authenticated, current link for each required step \\
R5 & source-scoped event log & initial-state, source-completeness, transition, order, and currentness premises \\
\bottomrule
\end{tabular}
\caption{The R1--R5 profile.}
\label{tab:r1r5}
\end{table}

R3 is deliberately disjunctive: a dependency may be bound in the declared
view, or a proof may show that varying it cannot affect the chosen observation.
This allows a portable abstraction layer without pretending the underlying
hardware or toolchain does not exist.

R4 follows the layout/link separation of \emph{in-toto}: a layout states the
required steps and authorized issuers, while a link carries occurrence-specific
attestation.  Authentication alone is insufficient.  The negative layout
accepts no issuer and therefore cannot be covered by an otherwise authentic
link.

\subsection{R5 and exact replay premises}

Let an event system have a step relation, an authoritative initial state, an
ordered authoritative event sequence, and a source-scoped retained log.  The
MeTTapedia replay theorem states
\[
  \mathsf{Trajectory}(s_0,E,s_f)
  \quad\Longleftrightarrow\quad
  \mathsf{replay}(s_0,L)=\mathsf{some}(s_f)
\]
only under all of the following premises:

\begin{enumerate}
  \item the retained and authoritative initial states agree;
  \item every event occurrence retains its required sources;
  \item retained events are complete as an occurrence multiset;
  \item the event order is valid and recovers the authoritative sequence;
  \item executable replay adequately implements the transition relation;
  \item the log receipt is current at the selected revision.
\end{enumerate}

Dropping a receipt reconstructs the wrong current state.  Reordering a complete
multiset can also fail replay.  A stale log may replay successfully yet fail
currentness.  These controls show why replayability, source completeness,
ordering, and currency remain distinct.

Prime DAG replay and HE derivation-local state-free execution instantiate this
boundary separately.  Their semantic objects are not collapsed into one
implementation-neutral trace format.

For a state-separating open-ended query schedule, a permanently sufficient
fixed context must be injective.  Therefore an evolving AGI may use bounded,
task-indexed views, but cannot recursively summarize away the authoritative
event source and still promise sufficiency for every future separating query.

\section{R6: protected recursive verifiability}

R6 is represented as an indexed certificate family.  The R1--R5 profile has
six fibres because R3 retains separate toolchain and hardware obligations.
A family-preserving modification supplies a map from every protected before
fibre to its corresponding after fibre.  A current modification additionally
carries an issued receipt whose revision matches the current revision.

\begin{theorem}[Restricted R6]
For a current family-preserving modification, an inhabited complete R1--R5
certificate family before the modification transports to an inhabited
complete family after the modification, and the modification receipt is
current.
\end{theorem}

The positive example is non-identity: it coarsens exact artifact observation to
an explicit input observation, transports each semantic certificate, and keeps
attestation and replay evidence.  The negative candidate preserves the build
relation and artifact observation but narrows the declared input view so that
hardware is hidden.  It cannot admit a protected-family map.  A second negative
uses a stale receipt.

This discharges Hatta's recursive-verifiability target for architectures whose
self-modifications carry the stated family-preservation, replay, and
currentness evidence.  It is not an unconditional theorem about arbitrary
self-modifying systems.

\section{Trusting trust and diverse double-compiling}

Ordinary recompilation can reproduce a malicious compiler executable from
apparently benign source when the executable recognizes both the login program
and its own source \cite{thompson1984}.  Matching the infected executable to
its self-recompiled output is therefore not sufficient evidence.

The formal DDC model distinguishes compiler source, compiler executables, the
trusted diverse compiler, source-language and executable-language semantics,
the first and second stages, deterministic equality at the selected
observation, and source correspondence.  Wheeler's proof assumptions are
represented as named fields rather than folded into a single ``trust''
predicate.

\begin{theorem}[Diverse double-compiling]
Under Wheeler's stated origin, compiler-correctness, determinism, trusted
diverse compiler, language-compatibility, and comparison assumptions, equality
of the DDC result with the compiler under test establishes correspondence to
the designated compiler source at the selected observation.
\end{theorem}

The development contains:

\begin{itemize}
  \item a substantive benign positive instance;
  \item a Thompson-style infected compiler that passes ordinary
        self-recompilation;
  \item a DDC experiment that detects the infected executable;
  \item for every proof-one and proof-two requirement, a countermodel satisfying
        all other requirements while the conclusion fails;
  \item a control showing that matching bytes do not themselves provide the
        proof-one premises;
  \item a hosting bridge that separates result equality from proof-fibre
        fidelity.
\end{itemize}

The C arithmetic differential example is used only as an observation-exact
hosting example.  It is not mislabeled as a DDC instance because it does not
construct Wheeler's complete two-stage compiler/source experiment.

\section{Observer-indexed verification legibility}

Legibility is not assigned to extensional behavior alone.  It is indexed by an
audit question, a view of the audit state, a presented realization, and an
auditor capability.

\begin{definition}[Audit information order]
An audit view $V$ is sufficient for question $q$ when $q$ factors through
$V$.  Write $V\preceq W$ when $W$ refines $V$.  Sufficient views are upward
closed in this order.  Two views are equivalent when each refines the other.
\end{definition}

Auditor capabilities enumerate permitted verification methods.  A
weakest-permitted method and an attained minimum rank must be supplied by
proof; the theory does not take an arbitrary inconvenient certificate as the
definition of opacity.

Hierarchical-complexity rank is a secondary readout.  Permutation-invariant
verification composition is classified as a chain; order-sensitive
composition is a coordination.  The same naturality square can be a chain at
endpoint observation and a coordination at proof-route observation.  Rank is
therefore observer-relative.

\begin{theorem}[Strict legibility separation]
There exist two presented realizations with extensionally equal behavior but
different attained minimum verification ranks.  The lower-rank realization
admits a direct evidence view; the higher-rank realization requires the full
ordered trace for its audit question.
\end{theorem}

Adding a higher-rank method cannot worsen legibility when a lower-rank method
remains permitted.  Conversely, low rank alone does not imply third-party
attestation or DDC feasibility: both require additional authentication and
compiler assumptions.  This prevents a scalar clarity score from becoming a
semantic authority.

\section{R7: coverage and resource feasibility}

R7 uses a verification plan with two independent readouts:

\begin{enumerate}
  \item the exact multiset of checked occurrences;
  \item a resource vector containing Work, Span, and retained-evidence storage.
\end{enumerate}

Sequential resource composition adds Work and Span.  Certified-independent
parallel composition adds Work, takes the maximum Span, and retains the sum of
evidence storage.  Both operations are associative and monotone; the parallel
readout is bounded by the sequential readout for the same components.

A verification requirement pairs a required case set with a resource limit.
Meeting the requirement means both covering the set and fitting the envelope.

\begin{theorem}[Sampling is scoped]
A verification certificate proves its property at every retained checked
occurrence.  It proves a declared scope only with a coverage proof, and proves
the universal statement only with full coverage.  A one-source sample of a
selectively deterministic build does not imply reproducibility at the omitted
source.
\end{theorem}

The same failure appears operationally for R5: a sampled reset receipt is
valid for that event, while omission of the later increment prevents full
trajectory coverage and reconstructs the wrong state.

Semantic reproducibility, current attestation, and economic feasibility are
pairwise separated by concrete controls.  A reproducible build can have an
audit plan exceeding its budget; a cheap plan can inspect an irreproducible
build; a current authentic attestation can coexist with either condition.

\section{Dynamic-link provenance and plural ledgers}

Hatta's protocol discussion asks composed tools to expose machine-readable
license and provenance metadata for downstream reasoning \cite{hatta2026}.
The formal route type decorates an existing proof-relevant semantic relation
with an ordered metadata history for every actual route witness.  Each link
records a link identity, tool identity, revision, source occurrences, and
license occurrences.

Relational composition retains the intermediate semantic witness and
concatenates the two histories.  Coverage of required link sets composes under
union.  Repeated source and license declarations remain repeated occurrences.

\begin{proposition}[Provenance boundaries]
The development proves:
\begin{enumerate}
  \item a compile--invoke chain covers both declared links;
  \item omitting the invocation ledger entry is detectable even when the
        endpoint relation is unchanged;
  \item extensionally equal routes may have different provenance histories;
  \item attaching metadata to a nondeterministic relation does not make it
        functional;
  \item compatible local ledgers glue only with an explicit gluing witness;
        mere runtime coordination does not manufacture one global ledger.
\end{enumerate}
\end{proposition}

Technical metadata and legal adjudication are modeled as separate fields.  A
factorization predicate states what would be required for artifact-plus-metadata
to determine adjudication.  Two cases with identical technical views and
different externally supplied adjudications refute that factorization in the
unconstrained model.  Thus metadata enables inspection but supplies no legal
rule by itself.  This is a logical non-entailment result, not a legal judgment.

\section{The theorem-level R1--R7 matrix}

The Lean development exports a matrix whose every row contains a formal
carrier, an actual positive proposition and proof, an actual negative
proposition and proof, an attribution tag, and its premise boundary.  Table
\ref{tab:coverage} summarizes the proof terms.

\begin{table}[ht]
\centering
\scriptsize
\begin{tabular}{@{}p{0.05\linewidth}p{0.30\linewidth}p{0.34\linewidth}p{0.21\linewidth}@{}}
\toprule
Req. & Positive & Negative & Boundary \\
\midrule
R1 & complete environment declaration & input-only view hides hardware & selected observation \\
R2 & exact identity build & exact branching build & selected budget \\
R3 & toolchain and hardware declared & hardware neither bound nor abstracted & declare-or-abstract \\
R4 & authorized current layout covered & rejecting layout uncovered & authorization, authentication, currency \\
R5 & complete log reconstructs state & lost receipt breaks completeness and replay & full replay premises \\
R6 & current non-identity coarsening transports family & behavior-only candidate and stale receipt fail & characterized modification class \\
R7 & full two-case audit meets envelope & sample is not universal; valid build can be unaffordable & coverage plus resources \\
\bottomrule
\end{tabular}
\caption{R1--R7 theorem coverage.}
\label{tab:coverage}
\end{table}

The matrix is not a checklist encoding: its top theorem eliminates the seven
cases and returns each row's concrete positive and negative proof terms.

\section{Trust boundary and limitations}

The mechanization contains no new axioms and no proof holes.  Crown theorems
depend only on standard Lean principles already used by the imported
foundations, chiefly quotient soundness, propositional extensionality, and
classical choice where representation selection requires it.

Several boundaries remain intentional.

\begin{itemize}
  \item Reproducibility at one observation does not upgrade itself to a finer
        observation.
  \item Result equality does not imply proof-fibre identity, provenance, or
        trusted attestation.
  \item R6 applies only to current family-preserving modifications.
  \item Sampled verification has no universal or probabilistic conclusion
        without an explicit coverage or probability theorem.
  \item Hierarchical rank is observer- and capability-indexed; it is not a
        universal human-readability measure.
  \item Dynamic-link metadata does not determine legal status.
  \item A concrete CeTTa/C-subset realization is downstream work.  Runtime
        executions may test or falsify a proposed observer, but do not define
        the source semantics they are meant to realize.
\end{itemize}

\section{Conclusion}

AGI-oriented reproducibility is most useful when treated as an invariant of a
declared observation and source view, not as a synonym for deterministic or
bit-identical execution.  Proof-relevant relational builds expose the exact
additional assumptions needed for functions, global ledgers, proof replay,
and third-party trust.  This makes Hatta's requirements compositional across
GSLT transformations and precise enough to survive restricted
self-modification, while the negative controls prevent provenance, sampling,
legibility, and cost from being mistaken for semantic authority.

\begin{thebibliography}{99}

\bibitem{hatta2026}
M.~Hatta.
\newblock Reproducibility is the new copyleft: Defining AGI-oriented
  reproducible builds.
\newblock In \emph{Artificial General Intelligence}, LNAI 16854, pages
  384--401, 2026.
\newblock doi:10.1007/978-3-032-33010-9\_25; arXiv:2606.03019.

\bibitem{hatta2025}
M.~Hatta.
\newblock Several issues regarding data governance in AGI.
\newblock arXiv:2508.12168, 2025.

\bibitem{thompson1984}
K.~Thompson.
\newblock Reflections on trusting trust.
\newblock \emph{Communications of the ACM}, 27(8):761--763, 1984.

\bibitem{wheeler2009}
D.~A. Wheeler.
\newblock \emph{Fully Countering Trusting Trust through Diverse
  Double-Compiling}.
\newblock PhD thesis, George Mason University, 2009.

\bibitem{lamb2021}
C.~Lamb and S.~Zacchiroli.
\newblock Reproducible builds: Increasing the integrity of software supply
  chains.
\newblock \emph{IEEE Software}, 2021. arXiv:2104.06020.

\bibitem{dolstra2006}
E.~Dolstra.
\newblock \emph{The Purely Functional Software Deployment Model}.
\newblock PhD thesis, Utrecht University, 2006.

\bibitem{dolstra2004}
E.~Dolstra, M.~de Jonge, and E.~Visser.
\newblock Nix: A safe and policy-free system for software deployment.
\newblock In \emph{LISA}, 2004.

\bibitem{mokhov2020}
A.~Mokhov, N.~Mitchell, and S.~Peyton Jones.
\newblock Build systems \`a la carte: Theory and practice.
\newblock \emph{Journal of Functional Programming}, 30:e11, 2020.

\bibitem{intoto2019}
S.~Torres-Arias, H.~Afzali, A.~Kuppusamy, R.~Curtmola, and J.~Cappos.
\newblock in-toto: Providing farm-to-table guarantees for bits and bytes.
\newblock In \emph{28th USENIX Security Symposium}, 2019.

\bibitem{malka2025}
G.~Malka, S.~Zacchiroli, and A.~Zimmermann.
\newblock Does functional package management enable reproducible builds at
  scale? Yes.
\newblock In \emph{Mining Software Repositories}, 2025.

\end{thebibliography}

\end{document}
